\subsection{Domination for extended measures}

To prove zero domination for possibly infinite measures, we construct a difference envelope with no mass where it is negative. This makes its entire measure finite and permits the preceding theorem to be applied. The resulting comparison yields domination for the original potentials.

\Needspace{10\baselineskip}
\begin{lemma}\label{lem:relative-zero-domination}
Let $u,v\in\mathcal E_\varphi$ satisfy $u\leq v$. If
\[
 \mu_u(\{u<v\})=0,
\]
then $u=v$.
\end{lemma}
\begin{proof}
Translate both potentials by a common constant so that $u\leq v\leq\varphi\leq0$. Fix $\varepsilon<\delta<1$, put $\beta=1-\delta$, and set
\[
 q=P_{\delta\theta}(v).
\]
This envelope is genuine by \cref{thm:relative-uniform-mass-lower}. By \cref{lem:relative-difference-envelope},
\[
 w=P_\theta(u-\varphi)\in\mathcal E_0,
 \qquad w+\varphi\leq u,\qquad w\leq0.
\]
The stretching lemma with reference zero shows that
$z=P_\theta(w/\beta)$ belongs to $\mathcal E_0$. Scaling the background and each defining competitor gives
\[
 \beta z\in\mathcal E_{\mathrm{env}}(X,\beta\theta,0).
\]
Indeed, if $a\in\PSH(X,\theta)$ lies below $t z$, then $\beta a\in\PSH(X,\beta\theta)$ lies below $t\beta z$, for every $t\geq1$.

Since $q\leq v\leq\varphi$, we have
\[
 \beta z\leq w\leq u-\varphi\leq u-q
\]
off the pole sets. Thus
\[
 k=P_{\beta\theta}(u-q)
\]
is genuine. Here the difference obstacle is taken off the pluripolar pole sets; it is quasi-continuous and finite almost everywhere. The general envelope theorem gives a finite upper bound for $k$ and the canonical inequality $k+q\leq u$. Upward stability therefore gives
$k\in\mathcal E_{\mathrm{env}}(X,\beta\theta,0)$.

Write $\nu_k=\mathrm{MA}^{\mathrm{ext}}_{\beta\theta}(k)$. Contact support carries $\nu_k$ by
\[
 D=\{k+q=u\}
\]
outside the pole sets. Positivity of mixed products and ordered contact give
\[
 \nu_k\leq\mathbf1_D\mu_u.
\]
These inequalities are first restricted to common finite truncation levels and then exhausted. On $D\cap\{k<0\}$ we have $u=k+q<q\leq v$, so $\nu_k(\{k<0\})=0$.

This vanishing proves that $\nu_k$ is finite. In fact, $\widehat k=\max(k,-1)$ is a single bounded $\beta\theta$-plurisubharmonic function, and locality gives
\[
 \nu_k(X)=\nu_k(\{k\geq0\})
 =\int_{\{k\geq0\}}\mathrm{MA}_{\beta\theta}(\widehat k)
 \leq\int_X\mathrm{MA}_{\beta\theta}(\widehat k)<\infty.
\]

Apply \cref{thm:relative-finite-domination} in the background $\beta\theta$, with reference zero and potentials $k,0$, using coefficient zero. Its hypotheses hold: $\beta\theta$ is semipositive and big, zero satisfies the model condition, and
\[
 v_-(\beta\theta)=\beta^n v_-(\theta)>0.
\]
The scaling identity follows by the bijection $a\mapsto\beta a$ between bounded potentials in the two backgrounds. We have just proved finiteness of $\nu_k$. The theorem yields $k\geq0$, hence $u\geq q$.

To pass to $\delta\uparrow1$, fix $\delta_0\in(\varepsilon,1)$ and $r=P_{\delta_0\theta}(v)$. For $\delta_0<\delta<1$, set $t=(\delta-\delta_0)/(1-\delta_0)$. The potential $t v+(1-t)r$ has background $\delta\theta$ and lies below $v$. Therefore
\[
 u\geq P_{\delta\theta}(v)\geq t v+(1-t)r.
\]
Let $t\uparrow1$ at points where $v,r$ are finite. This gives $u\geq v$ almost everywhere, and then everywhere by \cref{lem:qpsh-almost-everywhere-order}. Together with $u\leq v$, it proves equality.
\end{proof}

\begin{theorem}\label{thm:relative-extended-domination}\index{Domination principle!extended measures}
Let $u\in\mathcal E_\varphi$ and $v\in\PSH(X,\theta)$ with $v\preceq\varphi$. If $0\leq c<1$ and
\[
 \mu_u\leq c\mu_v\quad\hbox{on }\{u<v\},
\]
then $u\geq v$.
\end{theorem}
\begin{proof}
Put $w=\max(u,v)\in\mathcal E_\varphi$ and choose $b>1$ with $(1-b^{-1})^n>c$. The stretching lemma gives $p=Q_b(u,w)\in\mathcal E_\varphi$ and $p\leq u$. Its measure is carried by
\[
 D=\{b^{-1}p+(1-b^{-1})w=u\}.
\]
On $D$, positivity and ordered contact imply
\[
 b^{-n}\mu_p+(1-b^{-1})^n\mu_w\leq\mu_u.
\]
On $\{u<w\}$, locality gives $\mu_u\leq c\mu_w$. Restricting both inequalities to $D\cap\{u<w\}$ and to common finite truncation levels makes subtraction legitimate. Since $(1-b^{-1})^n>c$, it follows that $\mu_p$ vanishes there. Exhausting the finite levels proves that $\mu_p$ is carried by $D\cap\{u=w\}\subseteq\{p=u\}$. Thus $\mu_p(\{p<u\})=0$.

Apply \cref{lem:relative-zero-domination} to $p\leq u$, both in $\mathcal E_\varphi$. It yields $p=u$. Substitution into the defining convex inequality for $p$ gives $w\leq u$ off the poles, hence everywhere by \cref{lem:qpsh-almost-everywhere-order}. Thus $v\leq w\leq u$.
\end{proof}

\begin{corollary}\label{cor:relative-extended-exponential-order}\label{cor:relative-finite-exponential-order}
Under the hypotheses on $u,v$ in \cref{thm:relative-extended-domination}, for every $\lambda>0$,
\[
 \mathrm{e}^{-\lambda u}\mu_u\leq \mathrm{e}^{-\lambda v}\mu_v
 \quad\Longrightarrow\quad u\geq v.
\]
Also $\mu_u\leq c\mu_v$ on $X$ for a finite constant $c\geq0$ implies $c\geq1$.
\end{corollary}
\begin{proof}
On $\{u<v-a\}$ the exponential inequality gives $\mu_u\leq \mathrm{e}^{-\lambda a}\mu_v$. Apply the theorem to $v-a$ and let $a\downarrow0$. All weights are taken off the pluripolar pole sets. For the second assertion, if $c<1$, applying the theorem to $v+A$ for every real $A$ contradicts finiteness of both potentials at a common point.
\end{proof}

Applying the exponential implication in both directions proves uniqueness for a fixed positive exponent. Applied to two equations $\mu_u=c_1\nu$ and $\mu_v=c_2\nu$, the scalar assertion gives $c_1=c_2$. The final section shows why this does not determine the normalized potential at exponent zero.
